\documentclass{article}
\usepackage{ITB_MSCe}
\usepackage{tcolorbox}
\tcbuselibrary{breakable}
\usepackage{bold-extra} % bold tt
\usepackage{lipsum}
\usepackage{graphicx}
\usepackage{color}
\usepackage{amsthm}
\usepackage{amssymb}
\usepackage{amsmath}
\usepackage{amsfonts}
\usepackage{amstext}
\usepackage{latexsym}
\usepackage{upgreek}
\usepackage{hyperref}
\usepackage{url}
\usepackage{nicefrac}
\usepackage{pifont}
\usepackage{marvosym}
\usepackage{enumitem}
\usepackage{tikz}
\usepackage[numbers]{natbib}
\usepackage{eso-pic}
\usepackage{textcomp}
\usepackage{fancyvrb}
\usepackage{mwe}
\usepackage{lastpage}
\usepackage{listings, xcolor}
\usepackage{algorithmic}
\usepackage{algorithm}

\usepackage{array}
\usepackage{ragged2e}
%\bibpunct{}{}{;}{a}{,}{,}

\input{images/rgb}

\lstdefinestyle{python_code}{%
  language=Python, basicstyle=\ttfamily, columns=fullflexible,
  keywordstyle=\color{blue}, commentstyle=\itshape\color{green!50!black},
  stringstyle=\color{magenta}, showstringspaces=false,
  morekeywords={as,@staticmethod}, }

\lstset{ style=python_code }

\newsavebox{\sclbox}

\newcolumntype{L}[1]{>{\RaggedRight}p{#1}}



\newenvironment{declaration}
  {\noindent\textbf{Declaration:}} {\par}

\title{Explaining the policy--practice gap in AI education governance using
multi-layered NLP: A comparative study}

\author{
\normalsize
\begin{tabular}[t]{c@{\extracolsep{1em}}c} Mariem NSIRI & Geraldine GRAY\\
Technological University Dublin & Technological University Dublin\\
B00177462@myTUDublin.ie & supervisor.name@tudublin.ie
\end{tabular}
}

%\date{\today}

\begin{document}

\section{Data Preparation and Corpus Construction}

\subsection{PDF-to-Markdown Conversion}


The first preparation step was to convert the collected PDF documents into Markdown
files. This was done with two Docling-based pipelines: one for the policy corpus and
one for the European sentiment
corpus.\footnote{\url{https://github.com/nsirimai/mscdsa/tree/main/msc/progress/markdown_extraction}}
A direct PDF-to-CSV extraction approach, for example using tools such as
\texttt{pdfplumber}, would have been less suitable at this stage. Such tools can
extract text, but the result is often close to a flat text dump, where headings,
sections, figure captions, tables, page layout, and reading order are harder to
preserve. This would make later chunking, inspection, and interpretation more
difficult.

Markdown was therefore used as an intermediate format. It is easier to read, correct,
inspect, and transform than raw PDF output, while still preserving useful structure
such as titles, sections, lists, paragraphs, and some table or figure content.
Docling is a document-conversion toolkit that supports PDF understanding, layout
analysis, table structure recognition, and export to Markdown \cite{auer2024docling}.
Docling can also extract visual elements such as images and figures. Selected
figures, charts, screenshots, and visual summaries were then converted into textual
descriptions with ChatGPT and manually inserted into the Markdown files at the
relevant places. This helped preserve information that could be lost in a direct
PDF-to-CSV extraction. Markdown is therefore used here as a practical bridge between
raw PDFs and the structured textual datasets used for NLP.

A light Markdown cleaning step was then applied before creating the structured
datasets.\footnote{\url{https://github.com/nsirimai/mscdsa/tree/main/msc/progress/markdown_cleaning}}
This step did not rewrite or simplify the documents. It mainly removed technical
extraction noise, such as HTML comments, broken Markdown image links, raw local image
paths, isolated page numbers, table-of-contents artefacts, repeated headings,
decorative separators, and excessive blank spaces. Normal tables, headings, numbers,
section structure, and meaningful document content were preserved. The raw extracted
Markdown files and the cleaned Markdown files are available in the project repository
under the raw Markdown
folder\footnote{\url{https://github.com/nsirimai/mscdsa/tree/main/msc/markdown/raw}}
and cleaned Markdown
folder\footnote{\url{https://github.com/nsirimai/mscdsa/tree/main/msc/markdown/cleaned}},
respectively.


\subsection{Generation of Synthetic Sentiment Markdown Files}

Because the European sentiment corpus was smaller than the policy corpus, a synthetic
sentiment corpus was created as a support dataset for robustness
testing.\footnote{\url{https://github.com/nsirimai/mscdsa/tree/main/msc/markdown/synthetic/sentiment}}
The aim was not to create new evidence about public opinion, but to test whether the
NLP pipeline remains stable when the sentiment side is closer in size to the France
and Ireland policy corpus. This follows work showing that synthetic text can be
useful in limited-data settings, especially when the generated data is clearly
separated from the original corpus and used for controlled methodological purposes
rather than as a substitute for real observations
\cite{halterman2025synthetically,chen2023empirical}. The generation also reflects a
wider point in synthetic-data research: synthetic data may have practical utility,
but its value depends on how it is used, labelled, and evaluated
\cite{snoke2018general,chim2025evaluating}. A ChatGPT prompt was therefore used to
generate synthetic Markdown files in two simple styles: public-sentiment comments and
press-release-like
texts.\footnote{\url{https://github.com/nsirimai/mscdsa/tree/main/msc/prompts}} The
content focused on themes already present in the European sentiment corpus, including
AI literacy, teacher training, privacy, trust, transparency, student support,
responsible use, equity, and public confidence.

Several constraints were enforced to avoid contaminating the empirical data. The
synthetic files were not allowed to include numbers, percentages, fake survey
results, tables, or artificial quantitative claims. This was important because the
original sentiment reports contain real survey evidence, and invented figures could
be misleading if mixed with empirical results. The synthetic corpus was therefore
designed as qualitative auxiliary material, useful for corpus balance, sensitivity
checks, and pipeline stress testing. It was not designed to replace the original
sentiment evidence or to be fused with it as empirical data. This is consistent with
recent work on synthetic text evaluation, which argues that generated text should be
assessed through meaning preservation, divergence, fidelity, and downstream
usefulness before it is used in analysis \cite{chim2025evaluating,tari2026privacy}.

The main limitation is that the generated texts remain synthetic and can still be
distinguished from the original reports. This was confirmed through the validation
pipeline, especially the classifier two-sample test. This limitation is not treated
as a failure, but as a transparency signal. It confirms that the synthetic corpus
must remain labelled in the metadata and analytically separate from the original
corpus. The synthetic data is therefore used only in robustness analysis and
corpus-balance experiments. It is not used to make claims about real public opinion,
real survey frequencies, or actual respondents.



\subsection{Markdown-to-CSV Conversion}


After the Markdown files were cleaned and the synthetic sentiment Markdown files were
generated, the next step was to convert all Markdown documents into CSV
files.\footnote{\url{https://github.com/nsirimai/mscdsa/tree/main/msc/progress/csv_creator}}
This made the corpus easier to inspect and analyse with Python and NLP tools. This
step did not apply heavy cleaning or modify the document content. Its role was mainly
to organise each Markdown file into a structured row with metadata, including the
document identifier, corpus type, country, source type, file path, character count,
and word count. Keeping the full text in one CSV field made it easier to compare
document sizes, check the corpus structure, prepare later chunked datasets, and keep
a clear link between each text and its original Markdown file. The CSV files
therefore act as a structured bridge between the cleaned and synthetic Markdown
corpora and the later NLP pipelines. The raw CSV outputs are available in the project
repository.\footnote{\url{https://github.com/nsirimai/mscdsa/tree/main/msc/csvs/raw}}


\subsection{Data Chunking}

After the Markdown files were converted into CSV format, the next step was to create
a chunked version of the
corpus.\footnote{\url{https://github.com/nsirimai/mscdsa/tree/main/msc/progress/csv_chunker}}
The purpose of this step was not to clean or rewrite the documents, but to create
smaller and more comparable text units for NLP analysis. The policy reports and
sentiment reports are long and often contain several topics in the same document. If
a whole document is represented as one TF-IDF vector, embedding, or topic-modelling
input, different themes may be averaged together. This design is consistent with
earlier work on text segmentation, where long texts are divided into multi-paragraph
passages that correspond to subtopics \cite{hearst1997texttiling}. It is also
consistent with recent work on chunk embeddings, which notes that dense vector
retrieval can benefit from shorter text segments because long texts may be
over-compressed into a single embedding
\cite{gunther2024latechunking,duarte2024lumberchunker}. For this reason, the chunking
pipeline split the documents into heading-aware and size-controlled chunks, while
preserving the original document metadata. 

The chunker used Markdown headings as local context and stored this information in a
\texttt{heading\_context} field. It also preserved metadata such as \texttt{doc\_id},
corpus type, country, source type, file path, chunk index, character count, and word
count. The chunks were created around a target size of 350 words, with an overlap of
50 words and a minimum size threshold. The aim was to keep enough local context while
avoiding very long text units that could mix several themes. This is useful for later
embedding-based analysis, because Sentence-BERT represents texts as dense semantic
vectors that can be compared using cosine similarity \cite{reimers2019sentence}.
Applying embeddings to shorter chunks reduced the risk of merging unrelated themes
into one vector. The chunk-level comparison was therefore used as a practical design
choice to reduce over-compression of long, multi-topic documents into a single
representation. The generated chunked CSV files are available in the project
repository.\footnote{\url{https://github.com/nsirimai/mscdsa/tree/main/msc/csvs/chunked}}



The main strength of chunking is that it improves comparability between different
document types. Long policy reports, sentiment reports, and synthetic Markdown files
can be analysed using more similar units of text. It also improves traceability,
because each chunk remains linked to its source document and section context. This
supports a mixed computational and qualitative workflow: the model can operate on
manageable text units, while the researcher can still return to the original document
when interpretation is needed. This interpretive step remains important in
policy-text analysis, where topic modelling should support, not replace, qualitative
judgement \cite{isoaho2021topic}.

The main limitation is that chunking is not neutral. A chunk boundary can sometimes
split an argument, and overlap can repeat some text across neighbouring chunks. This
can slightly increase the chunk-level character count and may influence
frequency-based methods. G{\"u}nther et al. make a similar point in their work on
contextual chunk embeddings: standard chunking can lose information from the
surrounding context, so chunk size and overlap remain practical compromises rather
than perfect solutions \cite{gunther2024latechunking}. For this reason, the full
document-level CSV files were preserved alongside the chunked dataset.

\subsection{Synthetic Sentiment Distance Evaluation}

Before using the synthetic sentiment corpus in robustness experiments, its distance
from the original sentiment corpus was evaluated. The aim was not to prove that the
synthetic texts were identical to the original reports. Instead, the aim was to check
whether they were close enough to be useful as labelled auxiliary data, while still
keeping them separate from the empirical sentiment corpus. This follows recent work
on synthetic text evaluation, which argues that generated data should be assessed
through fidelity, divergence, and downstream usefulness rather than assumed to be
equivalent to real data \cite{chim2025evaluating,tari2026privacy}.

The validation used several complementary distance measures. Lexical similarity was
measured with unigram and bigram Jensen--Shannon divergence, which provides a
symmetric way to compare probability distributions \cite{endres2003new}. TF-IDF
centroid cosine similarity was also used to compare the vocabulary space of the
original and synthetic chunks, following the general vector-space model logic used in
information retrieval \cite{salton1988term}. Semantic similarity was measured with
sentence embeddings, using Sentence-BERT-style representations that can be compared
with cosine similarity \cite{reimers2019sentence}. To compare the broader embedding
distributions, the pipeline also used Fréchet embedding distance, inspired by
distributional evaluation in generative modelling \cite{heusel2017gans}, and Maximum
Mean Discrepancy, a kernel-based two-sample distance \cite{gretton2012kernel}.

A classifier two-sample test C2ST was also included. In this test, a classifier is
trained to distinguish original sentiment chunks from synthetic sentiment chunks. If
the classifier performs close to chance, the two samples are difficult to
distinguish; if it performs well, the synthetic corpus still contains detectable
differences. This follows the classifier two-sample testing framework proposed by
Lopez-Paz and Oquab \cite{lopez2016revisiting}. In this work, the C2ST result was
treated as a transparency signal. A high score does not automatically invalidate the
synthetic corpus, but it confirms that the synthetic data must remain labelled and
analytically separate from the original data.

The same pipeline also compared the European policy corpus with the original
sentiment corpus, and then with the balanced sentiment corpus made of original plus
selected synthetic sentiment chunks. This made it possible to check whether adding
synthetic sentiment data changed the policy--sentiment distance. The synthetic corpus
is therefore used only for robustness and sensitivity analysis. It is not used to
make claims about real public opinion, survey frequencies, or actual respondents.


\subsubsection*{Outcome of the Distance Evaluation}

The distance evaluation was mainly applied to the chunked
data.\footnote{\url{https://github.com/nsirimai/mscdsa/tree/main/msc/progress/synthetic_distance/chunk_based}}
The validation shows that the synthetic sentiment corpus reached a close
character-level balance with the France and Ireland policy corpus at chunk level. The
balanced sentiment corpus contains 1,323,166 characters, compared with 1,324,951
characters in the European policy corpus, leaving only a small difference of 1,785
characters. The original and synthetic sentiment chunks show moderate lexical
similarity, with a TF-IDF cosine score of 0.572, but stronger semantic similarity,
with an embedding centroid cosine of 0.951 and a low MMD value of 0.000391. This
suggests that the synthetic chunks cover a similar semantic space to the original
sentiment corpus, even if they do not reproduce the same surface wording.

The C2ST still shows high distinguishability, with an accuracy of 0.953 and a ROC-AUC
of 0.995. This confirms that the synthetic corpus is not identical to the original
corpus and should remain labelled and analytically separate. The European policy
comparison also shows that adding synthetic sentiment data does not artificially make
the sentiment corpus closer to policy language, since the TF-IDF cosine decreases
from 0.541 to 0.500. Overall, the outcome supports the intended use of the synthetic
sentiment corpus as auxiliary data for robustness and sensitivity testing, not as
empirical evidence about public opinion.

\subsubsection*{Comparison with Document-Level Validation}

A document-level validation was also run as a supplementary
check.\footnote{\url{https://github.com/nsirimai/mscdsa/tree/main/msc/progress/synthetic_distance/doc_based}}
At document level, the synthetic corpus showed weaker lexical and semantic alignment
with the original sentiment corpus than in the chunk-level pipeline. For example, the
TF-IDF cosine increased from 0.489 at document level to 0.572 at chunk level, while
the embedding centroid cosine increased from 0.877 to 0.951. The Fréchet embedding
distance also decreased from 0.821 to 0.342, and MMD decreased from 0.000757 to
0.000391. These results suggest that chunking gives a more faithful view of local
semantic similarity, because long documents often mix several themes in one
representation. However, both pipelines showed high classifier distinguishability,
with document-level ROC-AUC of 1.000 and chunk-level ROC-AUC of 0.995. This confirms
the same methodological conclusion: the synthetic sentiment corpus is useful for
robustness and balance testing, but it remains distinguishable from the original
corpus and must stay labelled and analytically separate.







\newpage
\bibliographystyle{abbrv}
\bibliography{bibliography}

\end{document}